% SPDX-License-Identifier: Apache-2.0
\section{Introduction}
\label{sec:intro}

\subsection{The evidence gap in PQC migration}

The migration to post-quantum cryptography is now a scheduling problem with a
regulatory deadline attached. NIST published ML-KEM~\cite{nist-fips203}, ML-DSA~\cite{nist-fips204} and
SLH-DSA~\cite{nist-fips205} in August 2024. On 12 October 2026 the German
Federal Ministry of the Interior (BMI), the Federal Ministry for Digital
Transformation and Government Modernisation (BMDS) and the Federal Office for
Information Security (BSI) published a joint national guideline implementing
the EU roadmap~\cite{bsi-pqc-migration}. It asks every organisation in scope to
finish an initial preparation and planning phase, including a cryptographic
inventory, by the end of 2026; to complete migration of high-risk use cases by
the end of 2030; and to have migrated nearly all systems by 2035
(Section~4.1 of the guideline). The dates are described there as urgent
recommendations rather than statutory deadlines, with the responsible bodies
asked to make them binding. Where this paper cites a section number of the
guideline, the number refers to the guideline and not to this paper; the two
numberings are unrelated.

The guideline is directly relevant to this paper for a second reason: it names
the artefact we build. Section~4.3 of the guideline recommends that each
organisation produce an overview of its cryptographic components, modelled on
the crypto register of the BSI IT-Grundschutz Compendium, and states that
``cryptographic inventory tools, repository scans and CBOMs (Cryptographic Bill
of Materials, e.g.\ based on the CycloneDX standard) can support the creation of
the cryptographic inventory.'' Section~4.6 of the guideline adds that providing
CBOMs simplifies comparison of cryptographic capabilities across the supply
chain, and that such provision ``can be set down in contracts.'' An inventory
artefact has thus moved from a good engineering practice to something a supplier
may be asked for contractually, which raises the cost of an inventory that does
not validate against its own standard.

Against this schedule, the binding constraint is not cryptographic primitives.
It is the ability to answer, with evidence, three questions: which cryptographic
assets are in a system, what has actually been proved about them, and what
remains open. An inventory that cannot be audited is a spreadsheet. A proof
whose gaps are not stated is a claim.

\subsection{Contributions}

This paper reports a case study along two axes, and one methodological finding
that emerged from combining them.

\begin{enumerate}[leftmargin=*]
  \item \textbf{A tiered Isabelle/HOL formalization of VWZ.} We mechanize the
        construction and organise 56 obligations into tiers by dependency
        (Section~\ref{sec:formalization}). Two tiers close with zero
        \texttt{sorry}; we state precisely which, and which do not.

  \item \textbf{A schema-validated CBOM in CI.} We generate a CycloneDX 1.6
        Cryptographic Bill of Materials for the implementation and gate it
        against the official schema on every push (Section~\ref{sec:cbom}).

  \item \textbf{Three documented failures of self-validation}, each caught only
        by an external reference: a build whose success means less than it
        appears to; a standard-format artefact that failed the standard while
        passing every internal test; and a user-facing page that never worked
        and was never noticed (Section~\ref{sec:case-study}).
\end{enumerate}

\subsection{What this paper does not claim}

We state this early because the failure mode we criticise is precisely the
overclaim.

We do not claim that VWZ is proved secure. Seventeen \texttt{sorry}s remain.
T3, the tier containing the security and probabilistic obligations, is entirely
open because its subjects are not yet defined. Key generation and the XOF are
not modelled at all. What is proved is a set of algebraic and linear-algebraic
facts about the construction, plus the correctness of a body of supporting
infrastructure.

We do not claim that a CBOM constitutes compliance with any regulation. It is an
inventory. Whether an inventory satisfies a regulator is a legal question.

We do not claim novelty in the cryptographic construction. The contribution is
engineering: how the verification was organised, what it cost, and what it
taught us about the reliability of our own evidence.

\subsection{Reproducibility}

Every number in this paper is re-derivable from the repositories. The
formalization ships with a script that recounts obligations and open goals from
the sources rather than from a table maintained by hand; the CBOM is regenerated
by a deterministic generator with a recorded schema hash; and the negative
control is a session that must fail.

The repository URLs, the snapshot commit identifiers and the Isabelle version
are given in Appendix~\ref{app:artefacts}.
